\section{The estimate for normal coordinates}\label{sec:estimate}
We work first with one bounded normal operator $Z$ and finitely many bounded selfadjoint coefficients $H_1,\ldots,H_m$.
Their represented star algebra is denoted by $\A$.
Assume that an outer derivation $\delta$ satisfies
\[
 \delta Z=1\otimes1,\qquad \delta H_l=0,\qquad
 \xi:=\delta^*(1\otimes1)\in H.
\]
The marginal law $\mu$ is diffuse, and $\delta Z^*=(\bar z-\bar w)/(z-w)$.
Restriction followed by projection to the marginal coarse space gives an actual one-coordinate derivation.
Thus the one-coordinate case of \cite[Proposition~3.5]{ACSY24} gives diffuseness, and the injective-gap argument of Lemma~\ref{lem:phase} determines the starred value in the full coarse space.
The formula for omitted coefficients is part of the present relative hypothesis.
Write $C:=\id\otimes\tau$ in this section.
Constants denoted by $c$ or $C_0,C_1,\ldots$ are absolute and may increase between estimates.

\subsection{Scalar control and liberation}
Let $\kappa(z,w,v)$ be the reciprocal circumradius of three distinct noncollinear planar points, with value zero for collinear or repeated triples.
Set
\begin{equation}\label{eq:scalar-budget}
 M_1\mu(z):=\sup_{r>0}\frac{\mu(B(z,r))}{r},\qquad
 G(\mu):=\int(M_1\mu)^2\dd\mu+\iiint\kappa^2\dd\mu^3.
\end{equation}

\begin{lemma}\label{lem:scalar}
If the scalar divided difference on $\C[Z,Z^*]$ has actual vacuum $\xi_\mu$, then
\begin{equation}\label{eq:scalar-estimates}
 G(\mu)\le c\|\xi_\mu\|_2^2,\qquad
 \left\|\int\frac{\overline{z-w}}{|z-w|^2+r^2}\dd\mu(w)\right\|_2
 \le\|\xi_\mu\|_2\quad(r>0).
\end{equation}
Its right slice is bounded from the uniform norm to $L^2(\mu)$ with norm at most $c\|\xi_\mu\|_2$.
\end{lemma}
\begin{proof}
Put $k_t(a):=\bar a/(|a|^2+t)$ and $C_t\mu(z):=\int k_t(z-w)\dd\mu(w)$.
The actual vacuum identity extends from polynomials to $C^1$ functions on a containing rectangle by polynomial approximation of the function and its first real derivatives.
The mean-value estimate controls the corresponding divided differences uniformly.

For a triple write $x:=|z-v|^2$, $y:=|w-v|^2$, $r:=|z-w|^2$ and $J:=2xy+2yr+2rx-x^2-y^2-r^2$.
Thus $J$ is sixteen times the squared triangle area.
The cyclic sum of $\Rea(k_t(z-w)\overline{k_t(z-v)})$ is
\[
 S_t=\frac{J+t(x+y+r)}{2(x+t)(y+t)(r+t)}\ge0.
\]
The cyclic sum of the real divided differences of $\overline{C_t\mu}$ has three-point kernel
\begin{equation}\label{eq:positive-soft}
 Q_t=\frac{J+\tfrac t2(x+y+r)+\tfrac{tJ}2(1/x+1/y+1/r)+3t^2}
 {(x+t)(y+t)(r+t)}\ge S_t.
\end{equation}
Both formulas follow by putting the three fractions over their common denominator; collision sets are product-null.
Fubini and the vacuum identity give
\[
 \|C_t\mu\|_2^2=\tfrac13\int S_t\dd\mu^3
 \le\tfrac13\int Q_t\dd\mu^3
 =\Rea\langle\xi_\mu,\overline{C_t\mu}\rangle.
\]
This proves the soft estimate, and Fatou gives $\int\kappa^2\dd\mu^3\le6\|\xi_\mu\|_2^2$.
For real $\|f\|_\infty\le1$, the positive measure $(1+f/2)\mu$ lies between $\mu/2$ and $3\mu/2$.
Positivity of $S_t$ therefore bounds its soft Cauchy transform in $L^2(\mu)$ by $\sqrt{27}\|\xi_\mu\|_2/2$.
Subtracting $C_t\mu$ and splitting complex $f$ into real and imaginary parts proves a uniform $L^\infty\to L^2$ bound for $C_t(f\mu)$.
For a scalar polynomial $p$,
\[
 pC_t\mu-C_t(p\mu)
 =\int\frac{p(z)-p(w)}{z-w}\frac{|z-w|^2}{|z-w|^2+t}\dd\mu(w)
 \longrightarrow E p
\]
in $L^2(\mu)$ by dominated convergence.
This proves the slice bound.

For completeness, the maximal-density term follows from a weighted commutator.
The little Grothendieck inequality \cite[Theorem~5.2]{Pis12} bounds the finite-column norm of $E:C(\operatorname{supp}\mu)\to L^2(\mu)$ by $c\|\xi_\mu\|_2$.
State factorization gives a state $\varphi$ and $c_0\le c\|\xi_\mu\|_2$ with $\|Ep\|_2^2\le c_0^2\varphi(|p|^2)$.
One may take a normal state with respect to $\mu$.
Indeed, the actual densely defined Hilbert adjoint of $E$ makes the bidual extension of $E$ vanish on the central summand killed by the representation on $L^2(\mu)$.
In the GNS factorization of $\varphi$, restriction to the remaining summand gives a normal positive functional of mass at most one with the same domination; adding the missing mass times $\mu$ gives a normal state.
The same bidual construction extends $E$ normally to $L^\infty(\mu)$, preserving $E1=0$, $[E,Z]=P_\Omega$ and the domination inequality.
Write it as $d\mu$ with $d\ge0$ and $\int d\dd\mu=1$.
The map $T(p\sqrt d):=Ep$ extends boundedly with norm at most $c_0$.
The commutator $[E,Z]=P_\Omega$ implies
\[
 |\tau(a)|\le2c_0\|Z-\lambda\|\|a\sqrt d\|_2.
\]
Testing at $(d+\varepsilon)^{-1}$ proves $d>0$ almost everywhere and $d^{-1}\in L^1(\mu)$.
Consequently $[T,Z]=|1\rangle\langle d^{-1/2}|$ on all of $L^2(\mu)$.
For a square $Q$ of side $\ell$, compress this identity and subtract its center from $Z$ to obtain
\[
 \mu(Q)^{1/2}\left(\int_Qd^{-1}\dd\mu\right)^{1/2}
 \le\sqrt2 c_0\ell,\qquad
 \frac{\mu(Q)}\ell\le\sqrt2 c_0\frac1{\mu(Q)}\int_Q\sqrt d\dd\mu.
\]
The second inequality uses Jensen's inequality for $s\mapsto s^{-2}$.
Four translated dyadic square filtrations contain every sufficiently small ball in a square of comparable side length.
Their $L^2$ maximal inequalities applied to $\sqrt d$ bound $\|M_1\mu\|_2$ by $cc_0$.
For radii larger than the support diameter $D$, use $1/D\le2c_0$, which follows from the uncompressed weighted commutator.
This completes \eqref{eq:scalar-estimates}.
\end{proof}

For a marginal tensor $a\otimes b$ let
\[
 I_{a\otimes b}(v\otimes w):=va\otimes bw.
\]
This is inner insertion and commutes with both outer actions.
The partial liberation derivation is
\[
 \gamma:=I_{z-w}\delta,
 \qquad \gamma f=f\otimes1-1\otimes f\quad(f\in\C[Z,Z^*]),
 \qquad \gamma H_l=0.
\]
Put $\zeta:=(i\gamma)^*(1\otimes1)=-i[Z^*,\xi]$.
It is selfadjoint.
Its left-right marginal spectral measure is
\[
 \nu(A\times B):=\|1_A(Z)\zeta1_B(Z)\|_2^2.
\]
This measure is symmetric, has zero diagonal mass, and satisfies
\begin{equation}\label{eq:gap-budget}
 I:=\int_{z\ne w}|z-w|^{-2}\dd\nu(z,w)\le\|\xi\|_2^2.
\end{equation}
Indeed $\zeta=-iS_{\bar z-\bar w}\xi$ for the left-right spectral calculus $S$.
The scalar vacuum is the orthogonal projection of $\xi$ onto $L^2(W^*(Z))$, so its norm is at most $\|\xi\|_2$.

For a bounded real symmetric marginal kernel $q$ and an analytic test
\[
 u=b\pm i\sqrt{1-b^2},\qquad b=b^*\in\A,\quad\|b\|<1,
\]
write
\[
 F_q(u):=\Rea\langle u,CI_q\gamma(u)\rangle.
\]
For a fixed test, this evaluation is continuous in the $L^2(\mu^2)$ norm of a kernel.
To see this, expand $\gamma b$ into finitely many bounded outer coefficients and differentiate the binomial series; its derivative terms are bounded by a summable multiple of $n\|b\|^{n-1}$.

\begin{lemma}\label{lem:kernel-energy}
For a marginal projection $e$ of mass $t>0$,
\begin{equation}\label{eq:projection-square}
 tF_{e\otimes e}(u)=\|iCI_{e\otimes e}\gamma(u)-ue\zeta e\|_2^2-\|e\zeta e\|_2^2.
\end{equation}
For a nonzero bounded real marginal function $k$,
\begin{equation}\label{eq:uncentered}
 F_{k\otimes k}(u)\le(\tau k)^2+\frac{\|k\zeta k\|_2^2}{\|k\|_2^2}.
\end{equation}
If $q$ is bounded, real, symmetric and supported on $B^2$, where $t:=\mu(B)$, then
\begin{equation}\label{eq:localized}
 |F_q(u)|\le\|q\|_\infty\bigl(2\sqrt{t\nu(B^2)}+3t^2\bigr).
\end{equation}
\end{lemma}
\begin{proof}
The marginal tensor adjoint formula is
\[
 (i\gamma)^*q=S_q\zeta+i(\operatorname{col}q-\operatorname{row}q).
\]
The row and column maps are bounded on $L^2(\mu^2)$.
Closedness of this adjoint implies $\nu\ll\mu^2$: approximate an indicator of a $\mu^2$-null set in $L^2(\mu^2+\nu)$ by polynomial tensors and apply the formula.
The graph closure of the marginal restriction has domain $L^2(\mu^2)\cap L^2(\nu)$ and the same formula.
In particular $\beta_q:=iI_q\gamma$ has actual vacuum $\eta_q=S_q\zeta$ for real symmetric bounded $q$.

Here are the finite-kernel calculations needed below.
Put $E_q:=C\beta_q$, $a_q:=\operatorname{row}q$ and $q^{\circ2}(z,v):=\int q(z,w)q(w,v)\dd\mu(w)$.
Put $C_q:=(\id\otimes\tau\otimes\tau)((\beta_q\otimes1)\beta_q-(1\otimes\beta_q)\beta_q)$.
The product rule gives
\begin{align*}
 C_qp&=i\{E_q(p)a_q+E_q(pa_q)-pE_q(a_q)-E_{q^{\circ2}}p\},\\
 \Rea\langle u,C_qu\rangle
 &=-2\Ima\tau(a_qu^*E_qu)+\Ima\langle u,E_{q^{\circ2}}u\rangle,\\
 \|E_qu\|_2^2&=2\Rea\langle E_qu,u\eta_q\rangle-\Rea\langle u,C_qu\rangle.
\end{align*}
One obtains the last identity first on polynomials, with the additional term $-\Rea\langle\beta_q(p^*p),1\otimes\eta_q\rangle$, and then passes to $u$ by the differentiated binomial series.
This passage uses first derivatives of the defect, which is itself an outer derivation, and requires no second-iterate domain assertion at $u$.
Bounded Borel marginal factors are allowed by uniformly bounded strong-star graph approximation in the liberation derivation.

For $q=e\otimes e$, use $a_q=te$, $q^{\circ2}=tq$ and the right support $E_qu=E_qu e$ to obtain \eqref{eq:projection-square}.
For $q=k\otimes k$ the same identities give
\[
 \|k\|_2^2F_q(u)=\|k\zeta k\|_2^2-\|E_qu-uk\zeta k\|_2^2
 +2(\tau k)\Ima\tau(ku^*E_qu).
\]
Completing the square against $i(\tau k)uk$ proves \eqref{eq:uncentered}.
For \eqref{eq:localized}, normalize $\|q\|_\infty\le1$ and put $q_j:=q^{\circ j}$.
Then $\|q_j\|_\infty\le t^{j-1}$, $\|\eta_{q_j}\|_2\le t^{j-1}\nu(B^2)^{1/2}$ and $\|a_{q_j}\|_2\le t^{j+1/2}$.
The displayed energy identities and right support give, with $v_j:=\|E_{q_j}u\|_2$,
\[
 v_j^2\le2(t^{j-1}\nu(B^2)^{1/2}+t^{j+1/2})v_j+\sqrt t\,v_{2j}.
\]
For this fixed test the finite expansion gives $v_j\le C_ut^j$.
Thus the finite supremum of $v_j/t^{j-1}$ over dyadic $j$ is at most $2\nu(B^2)^{1/2}+3t^{3/2}$.
Pairing once more with $u1_B$, of norm $\sqrt t$, proves \eqref{eq:localized}.
Conditional expectations onto finite product partitions preserve symmetry, support and the uniform norm and converge in $L^2(\mu^2)$, so the conclusion passes from finite kernels to bounded measurable ones.
\end{proof}

\subsection{The mixed curvature energy}
Put $\delta^\#p:=(\delta(p^*))^\dagger$, where $(v\otimes w)^\dagger:=w^*\otimes v^*$, and set
\[
 \delta_x:=(\delta+\delta^\#)/2,\qquad
 \delta_y:=(\delta-\delta^\#)/(2i),\qquad T_a:=C\delta_a\quad(a=x,y).
\]
Their actual selfadjoint vacua are $\zeta_x=(\xi+\xi^*)/2$ and $\zeta_y=i(\xi-\xi^*)/2$, with squared norms summing to $\|\xi\|_2^2$.
For
\[
 g(z,w):=\int\kappa(z,v,w)^2\dd\mu(v),\qquad
 V(u)^2:=\|T_xu\|_2^2+\|T_yu\|_2^2,
\]
the required energy identity is
\begin{equation}\label{eq:mixed-energy}
 V(u)^2=2\Rea\sum_{a=x,y}\langle T_au,u\zeta_a\rangle-\tfrac12F_g(u).
\end{equation}
The kernel $g$ need not belong to $L^2(\mu^2)$.
Its use here is through the generator tensors $(z-w)g$ and $(\bar z-\bar w)g$.
They do belong to $L^2(\mu^2)$, since Jensen's inequality and $|z-w|\kappa(z,v,w)\le2$ give
\begin{equation}\label{eq:gap-g}
 \int|z-w|^2g(z,w)^2\dd\mu^2\le4\int\kappa^2\dd\mu^3.
\end{equation}

We give the domain justification for \eqref{eq:mixed-energy}.
The smooth phases $h_t(z,w):=(\bar z-\bar w)^2/(|z-w|^2+t)$ have modulus at most one.
Their first-leg divided differences, the second-leg divided differences and their conjugate-denominator versions have squared modulus at most $8Q_t$ from \eqref{eq:positive-soft}.
For example, writing $a=z-v$, $b=w-v$, the numerator identity
\[
 \bar a^2(|b|^2+t)-\bar b^2(|a|^2+t)
 =\bar a\bar b(\bar a b-\bar b a)+t(\bar a-\bar b)(\bar a+\bar b)
\]
gives this bound by separating the area term and the remaining positive terms of $Q_t$.
Its integral is at most $24\|\xi_\mu\|_2^2$ by the scalar vacuum calculation.
Approximate $h_t$ in $C^1$ on a containing real box by polynomial tensors.
Weak closedness of the joint graph of the four tensor-leg derivatives puts both phases in that original joint graph closure as $t\downarrow0$.
Convex combinations turn the weak graph convergence into strong graph convergence while keeping uniform bounds for the tensors.
Their limiting first divided difference is
\[
 \frac{(\bar z-\bar v)(w-v)-(\bar w-\bar v)(z-v)}
 {(z-v)(w-v)(z-w)},
\]
whose squared modulus is $\kappa(z,w,v)^2$.
The same approximants work in the mixed tensor-leg graphs.
Multiplication by each fixed core outer coefficient is continuous in these graphs; the Leibniz rule therefore supplies actual joint graph approximants to $\delta p$ and $\delta^\#p$ for every $p\in\A$.
Their partial traces are uniformly operator-norm bounded and converge in the first-derivative graphs, so their products with each fixed $L^2$ vacuum converge in $H$.
The actual adjoint formula for each real slice is $T_a^*p=p\zeta_a-T_ap$.

Let $A_a:=(\delta_a\otimes1)\delta_a-(1\otimes\delta_a)\delta_a$.
The cross terms in its Leibniz rule cancel.
The just established graphs justify contracting these actual iterates and pairing with a polynomial.
For the sum $A_x+A_y$, direct evaluation at $Z,Z^*$ gives, after tracing the middle leg, the tensors $(z-w)g/2$ and $(\bar z-\bar w)g/2$; all omitted coefficients are annihilated.
For instance, if $N:=\bar a b-\bar b a$, the relevant numerator at $Z$ is $-N^2/(\bar a|b|^2|a-b|^2)=a\kappa(a,b,0)^2$ before the factor $1/2$.
Equality on generators and the outer Leibniz rule identify the whole defect with one half of the weighted liberation.
The real-slice energy calculation used in Lemma~\ref{lem:kernel-energy} now proves \eqref{eq:mixed-energy} on analytic tests.
Positive kernel approximations dominated by $g$ pass to this identity by \eqref{eq:gap-g}, the finite word rule and the differentiated binomial series.
Consequently a bound $F_g(u)\ge-c\|\xi\|_2^2$ implies $V(u)\le c\|\xi\|_2$.
Averaging the two signs in $u=b\pm i\sqrt{1-b^2}$, then scaling and splitting real and imaginary parts, gives the right-slice estimate on $\A$.

\subsection{Disks with small local mass}
Fix a smooth decreasing cutoff $\psi$ equal to one on $[0,1]$ and zero on $[2,\infty)$, and define
\[
 \theta(x,r):=r^{-1}\int\psi(|x-y|/r)\dd\mu(y).
\]
We need the following square estimate without a uniform growth assumption on $\mu$:
\begin{equation}\label{eq:density-square}
 \int\int_0^\infty|\theta(x,r)-\theta(x,2r)|^2\frac{dr}{r}\dd\mu(x)
 \le c_\psi G(\mu).
\end{equation}
Here is a reduction to three established harmonic-analysis results, with their hypotheses made explicit.
Write $b(x)^2:=\iint\kappa(x,y,z)^2\dd\mu(y)\dd\mu(z)$, $f:=M_1\mu+b$ and $\sigma:=\mu/f$ on the full-measure set where $f$ is finite.
The positive lower bound for $M_1\mu$ on its compact support makes $\sigma$ finite.
For every disk, $\sigma(B(a,r))\le2r$, since $M_1\mu(x)\ge\mu(B(a,r))/(2r)$ for $x\in B(a,r)$.
For every Borel set $Q$,
\[
 \iiint_{Q^3}\kappa^2\dd\sigma^3\le3\sigma(Q),\qquad
 \int f^3\dd\sigma=\int f^2\dd\mu\le2G(\mu).
\]
For the first estimate, assign each triple to a point with smallest $f$ and bound the reciprocal product by that point's $f^{-3}$.

Tolsa's global density-square estimate under linear growth \cite[Theorem~19.1, arXiv version]{Tol17}, applied to $\sigma|_Q$, bounds its sharp square integral by $c\sigma(Q)$.
Smooth density differences are positive averages of sharp ones with weights $s(-\psi'(s))\,ds$, $1\le s\le2$, so the same testing bound holds for the smooth operator $U$ with values in $\mathcal K:=L^2((0,\infty),dr/r)$.
Its scalar kernel at scale $r$ is supported on $|x-y|\le4r$, is bounded by $c/r$, and has both first gradients bounded by $c/r^2$.
The indicator tests $b_Q=1_Q$, the linear growth of $\sigma$ and these size and second-variable regularity estimates verify \cite[Theorem~1.3]{LM17} with $m=\alpha=1$.
Thus $U:L^2(\sigma)\to L^2(\sigma;\mathcal K)$ is bounded uniformly.
The $\mathcal K$-valued kernel has size $c/|x-y|$ and difference bound $c|x-x'|/|x-y|^2$ when $|x-x'|\le|x-y|/2$, in both variables.
Linear growth and annular summation give both integral H\"ormander conditions; the component formula represents $U$ away from the support of the input.
These are precisely the additional hypotheses in \cite[Definition~2.1 and Theorem~2.7]{GM00}, which yield the strong $L^3(\sigma)\to L^3(\sigma;\mathcal K)$ bound.
Since $\mu=f\sigma$, H\"older's inequality gives
\[
 \int\|Uf\|_{\mathcal K}^2f\dd\sigma
 \le\|Uf\|_{L^3(\sigma;\mathcal K)}^2\|f\|_{L^3(\sigma)}
 \le c_\psi\int f^3\dd\sigma.
\]
This proves \eqref{eq:density-square} for the original measure.

For a disk $B=B(a,r)$ put
\[
 T(B):=\int\min(1,r^2/|a-v|^2)\dd\mu(v),\qquad
 d\lambda(B):=dA(a)\,dr/r^5,
\]
where $dA$ is planar area measure.
Fix an absolute large $K$ and call a disk small in local mass when $T(B)>K\mu(3B)$.
For $x\in B$, annular decomposition and $r\theta(x,r)\le\mu(3B)$ show that, on these disks,
\[
 \frac{T(B)^2}{r^2}
 \le c\sum_{j\ge0}2^{-j}|\theta(x,2^jr)-\theta(x,2^{j+1}r)|^2.
\]
Indeed $T(B)/r\le c_0\sum_j2^{-j}\theta(x,2^jr)$; for $K>4c_0$ the initial constant term is absorbed, after which telescoping and weighted Cauchy--Schwarz give the assertion.
Integrating over $x\in B$ and the centers of the disks gives, with $\mathcal B$ the selected disks,
\begin{equation}\label{eq:small-budgets}
 L:=\int_{\mathcal B}\mu(B)T(B)^2\dd\lambda(B)\le cG(\mu),\qquad
 N:=\int\nu(B^2)\dd\lambda(B)\le cI.
\end{equation}
The first bound uses \eqref{eq:density-square} and the area $\pi r^2$ of possible centers for a fixed $x$.
For the second, a disk containing $z,w$ has $r\ge|z-w|/2$ and the same center-area bound; integration in $r$ costs $c|z-w|^{-2}$.

Set $h_v(z):=(z-v)/(\bar z-\bar v)$ away from $v$.
Its squared oscillation on $3B$, integrated in $v$, is at most $cT(B)$, by differentiation when $|a-v|>6r$ and the bound $|h_v|\le1$ otherwise.
Thus the collected angular kernel
\[
 q_B(z,w):=1_B(z)1_B(w)\int|h_v(z)-h_v(w)|^2\dd\mu(v)
\]
satisfies $\|q_B\|_\infty\le cT(B)$.
By \eqref{eq:localized} its absolute signed cost is at most
\[
 cT(B)\{\sqrt{\mu(B)\nu(B^2)}+\mu(B)^2\}.
\]
As $\mu(B)\le T(B)$, its integral over $\mathcal B$ is bounded by $c(\sqrt{LN}+L)\le c\|\xi\|_2^2$.

\subsection{The signed estimate on a larger disk}
Assume temporarily that a finite marginal partition $(C_j)$ satisfies $\nu(C_j^2)=0$ for every $j$.
We prove that for $B=B(a,r)$, $Q=3B$, $m=\mu(Q)>0$ and real bounded $p$ on $Q$,
\begin{align}
 F_{1_B(z)1_B(w)(p(z)-p(w))^2}(u)
 \ge{}&-\int_{B\times(Q\setminus B)}|p(z)-p(w)|^2\dd\mu(z)\dd\mu(w)\notag\\
 &-\frac{8\osc_Q(p)^2}{m}\nu(B\times Q).\label{eq:disk}
\end{align}
Zero-mass disks contribute zero.
First let $p$ be simple and partition $Q$ into positive-mass cells refining $(C_j)$, $B,Q\setminus B$ and the level sets of $p$.
Write $t_i$ for their masses, $p_i$ for their values, $s_i=1$ on $B$ and zero otherwise, and $\nu_{il}$ for the corresponding spectral rectangles.
Let $F=(f_{ij})$ be the real symmetric matrix representing the signed pairings of cell kernels, with the off-diagonal entries symmetrized.
The projection square identity gives $f_{ii}\ge0$, also for any subcell, since $\nu_{ii}=0$.
Put $t=(t_i)$ and $X_0:=tt^{\mathsf T}-F$.
For every pivot $l$,
\begin{equation}\label{eq:row-positive}
 X_0+D_l\succeq0,\qquad D_l:=\operatorname{diag}(2\nu_{il}/t_l)_i.
\end{equation}
To check this, split cell $l$ into $N$ sets $R$ of mass $t_l/N$ and apply \eqref{eq:uncentered} at $k\pm N^{3/4}1_R$ for a fixed real cell function $k$.
Average the signs and discard the extra nonnegative term $N^{3/2}F_{1_R\otimes1_R}$.
The squared denominator is $t_l\sqrt N(1+o(1))$, uniformly over $R$, while the trace perturbation tends to zero.
Since $\nu(R^2)=0$, the averaged numerator is exactly
\[
 \|k\zeta k\|_2^2+N^{3/2}\int\{k(z)^21_R(w)+1_R(z)k(w)^2\}\dd\nu.
\]
Averaging the $N$ choices of $R$ and using symmetry of $\nu$ yields $2\sum_i k_i^2\nu_{il}/t_l$ in the limit.
The normalized positive numerators have bounded average, so the uniform relative denominator error may be removed.
This proves \eqref{eq:row-positive}.

Center $p$ in the $Q$-weights, so $\sum_it_ip_i=0$, and put $v_0:=m^{-1}\sum_it_ip_i^2$.
On all cells of $Q$ form the weighted Laplacian with off-diagonal entries $-(p_i-p_j)^2$ and diagonal $\sum_j(p_i-p_j)^2t_j/t_i$.
Compress it by $\operatorname{diag}(s_i)$ and call the resulting positive matrix $L$.
Then
\[
 t^{\mathsf T}Lt=\sum_{i\in B,\ j\in Q\setminus B}t_it_j(p_i-p_j)^2.
\]
There is a useful positive decomposition $L=\sum_lA_l$, where
\begin{align*}
 r_l(i)&:=\sqrt{t_l/m}\sqrt{p_l^2+v_0}\,s_i((m/t_l)1_{i=l}-1),\\
 b_l(i)&:=\sqrt{2t_l/m}\,s_ip_i,\qquad
 A_l:=r_lr_l^{\mathsf T}+b_lb_l^{\mathsf T}.
\end{align*}
The off-diagonal sum is $-s_is_j(p_i-p_j)^2$.
Before compression each vector is perpendicular to $t$, which determines the diagonal of the sum uniquely.
For $i\ne l$,
\[
 (A_l)_{ii}\le4(t_l/m)s_i\osc_Q(p)^2.
\]
No bound is needed at $i=l$, since $(D_l)_{ll}=0$.
With $W_{ij}:=s_is_j(p_i-p_j)^2$, the matrix $W+L$ is nonnegative and diagonal, whence
\[
 \operatorname{tr}(WF)\ge-\operatorname{tr}(LF)
 =\operatorname{tr}(LX_0)-t^{\mathsf T}Lt
 \ge-\sum_l\operatorname{tr}(A_lD_l)-t^{\mathsf T}Lt.
\]
This proves \eqref{eq:disk} for simple functions.
Uniform quantization of bounded real functions and the fixed-test insertion continuity extend it to bounded Borel $p$.

The disk identity
\[
 \int1_B(z)1_B(w)\dd\lambda(B)=\frac{c_*}{|z-w|^2}\quad(z\ne w)
\]
has $0<c_*<\infty$ by translation and dilation.
Also $|h_v(z)-h_v(w)|^2=|z-w|^2\kappa(z,w,v)^2$ off coincidences.
Apply \eqref{eq:disk} to the real and imaginary parts of $h_v$ and integrate against $c_*^{-1}d\mu(v)d\lambda(B)$.
For its first term, $z\in B,w\in3B$ forces $r\ge|z-w|/4$ and a center area at most $\pi r^2$.
The cost is at most $c\int\kappa^2\dd\mu^3$.
On disks outside $\mathcal B$, the integrated second term costs at most
\[
 cK\int\nu(B\times3B)\dd\lambda(B)\le cKI.
\]
On $\mathcal B$, use the absolute bound already obtained from \eqref{eq:small-budgets}.
At bounded radii these are ordinary scalar Fubini computations for each test.
The positive kernels exhaust $g$, and \eqref{eq:gap-g} passes their signed evaluations to the limit.
We have proved
\begin{equation}\label{eq:partition-bound}
 F_g(u)\ge-c\|\xi\|_2^2,\qquad
 \|C\delta p\|_2\le c\|\xi\|_2\|p\|,
\end{equation}
under the auxiliary finite-partition hypothesis, with a constant independent of its masses and cardinality.

\subsection{Passage to the original tuple}
The auxiliary partition cannot be imposed on an arbitrary original coordinate.
Lemma~\ref{lem:lattice-transfer}, proved in Appendix~\ref{sec:reductions}, transfers any finite lattice model to an actual relative coordinate on finitely many separated circles, with zero liberation mass on every individual circle square.
The relative vacuum costs at most a fixed multiple of the lattice budget, and conditional expectation recovers the lattice derivative on every old word.
Thus \eqref{eq:partition-bound} gives a uniform lattice estimate.
The reverse transfer, by rounding a diffuse spectral law, gives the same uniform bound for every relative model with one trace-half projection coefficient.
Lemma~\ref{lem:projection-transfer} then passes the bound, with at most a factor two, to finitely many selfadjoint coefficients.

For a normal generating tuple with $\mathsf A$, Lemma~\ref{lem:phase} says that $\partial_i$ kills the adjoints of the omitted coordinates.
Their real and imaginary parts are therefore exactly coefficients of the permitted kind.
We obtain the right-slice estimate of \eqref{eq:main-slices} on the original $\A_X$.
The derivation $p\mapsto(\partial_i(p^*))^\dagger$ is the normalized coordinate derivation for the adjoint tuple $X^*$, with actual vacuum $\xi_i^*$.
Its right slice is the tracial involution of the original left slice at $p^*$.
Applying the same estimate gives the other slice with the same absolute constant.
Lemma~\ref{lem:graph} restores both actual domains of $\overline{E_i^\circ}$ and its adjoint and passes the commutator to $\M\Omega$.
Negating these closures proves $\mathsf A\Rightarrow\mathsf S$ and completes Theorem~\ref{thm:normal}.
